\chapter{References, Code and Appendices}\label{cha:refs}

This chapter continues with citing the literature, code, and appendices. The examples of the sections do not depend on one another; copy whichever you need and replace it with your own thesis data.

The files for the references and appendices live in \code{contents/back/}, and the other examples can be written straight into this chapter. When moving or editing content, remember to update the matching lines in \code{main.tex} at the same time, so that the order of the references or appendices does not slip.

\section{Citing the Literature}

A \code{.bib} file is a plain-text bibliography database in \BibTeX{} format; one entry corresponds to one work, in the form \code{@type{key, field = {value}, ...}}. \code{contents/back/references.bib} already holds a few sample entries, one of which looks like this:

\begin{lstlisting}
@inproceedings{kocher99,
  author =       {C. Kocher and J. Jaffe and B. Jun},
  title =        {Differential Power Analysis},
  editor =       {M. Wiener},
  booktitle =    {Advances in Cryptology ({CRYPTO}~'99)},
  series =       {Lecture Notes in Computer Science},
  volume =       1666,
  pages =        {388-397},
  publisher =    {Springer-Verlag},
  month =        {August},
  year =         1999
}
\end{lstlisting}

\code{@inproceedings} is the entry type (a conference paper); others you will often see are \code{@article} (journal article), \code{@book} and \code{@incollection} (a chapter in a book), and the type affects which fields the reference list shows. \code{kocher99} is the key used to cite the entry in the text and must not repeat within one \code{.bib}; the \code{author}, \code{title}, \code{year} and other fields that follow are the data of the work, filled in one by one, and their order does not affect the result.

Below each search result, Google Scholar has a ``Cite'' button, and at the bottom of its pop-up window is a \BibTeX{} link that copies the whole entry; the pages of library databases (such as IEEE Xplore and the ACM Digital Library) mostly have an ``Export Citation'' or ``Cite'' function too, where you choose the \BibTeX{} format. If you usually work with a reference manager such as Zotero or EndNote, you can also export a \code{.bib} file directly. Once you have an entry, paste it into \code{references.bib}, change the key to a name you can remember (for example \code{surname+year}), and check that it does not repeat an existing entry.

With the data ready, go back to the chapter file and cite with the same key. To put a work at the end of a sentence, use \code{\citep}; to put the author into the sentence, use \code{\citet}. The forms that show only the author or the year are listed too, so choose whichever the sentence needs. These four commands come from the \code{citation-style-language} package, which the template already loads.

\begin{lstlisting}[language=latex]
\citep{kocher99}       % parenthetical citation
\citet{kocher99}       % in-text citation
\citeauthor{kocher99}  % author only
\citeyear{kocher99}    % year only
\end{lstlisting}

The differences among the four commands after compiling are as follows; check them first and then replace the key with the work you want to cite:

\begin{tlt}{Comparison of citation commands}{tbl:cite-forms}{L{0.3\linewidth}L{0.25\linewidth}L{0.35\linewidth}}
Command & Result & Use \\
\midrule
\code{\citep{kocher99}} & \citep{kocher99} & Parenthetical, at the end of a sentence \\
\code{\citet{kocher99}} & \citet{kocher99} & In-text, the author is part of the sentence \\
\code{\citeauthor{kocher99}} & \citeauthor{kocher99} & Shows only the author \\
\code{\citeyear{kocher99}} & \citeyear{kocher99} & Shows only the year \\
\end{tlt}

To cite a particular passage in a book, add a page or chapter in square brackets. The example below uses the entry \code{companion} and adds ``page 42''; when you apply it to your own thesis, change it to the location you actually cite:

\begin{lstlisting}[language=latex]
\citep[p.~42]{companion}
\end{lstlisting}

After compiling it appears as \citep[p.~42]{companion}. The page is typeset together with the citation, and the punctuation and arrangement are decided by the reference style you chose.

To cite several works at once, separate the keys with commas, for example \code{\citep{kocher99,IEEE-1363}}. Like a single citation, it uses the reference style from the config, so pick whichever way you are more used to.

\section{Reference Settings}

Open \code{config.yml} and add the settings you need under the existing \code{bibliography:} block. \code{path} names the bibliography files and \code{style} decides the citation format; if your data are spread over several files, list the paths in \code{path} in order.

\begin{lstlisting}[language=yaml]
bibliography:
  path: ["/contents/back/references.bib"] # several files allowed
  style: "apa"       # built-in CSL style name, or the path of a .csl file; see the table below
  title: "參~考~文~獻"
  color: "#0000ff"   # colour of citations and links; use "#000000" for black-and-white print
  full: false        # true also lists entries that were not cited
\end{lstlisting}

\code{path} may mix \code{.bib}, Hayagriva \code{.yml} (the format of the Typst version) and CSL JSON/YAML; Hayagriva files are converted to CSL JSON automatically when \code{scripts/gen.lua} runs, so you need not convert them yourself. A citation key cannot appear in two different files.

\code{style} decides the logic of the references and falls roughly into three systems: author-date shows the author's name and year in the text, so readers see at a glance how old the study is; author-page shows the author and page, which makes it easy to check the source; and numeric keeps only a number or superscript in the text and orders the list by number, the least distracting to read. Those in common use in each field are:

\begin{tlt}[notes={\texttt{mla} is short for \texttt{modern-language-association}; either spelling works.}]{Common reference formats}{tbl:citation-styles}{L{0.12\linewidth}L{0.17\linewidth}L{0.25\linewidth}L{0.26\linewidth}}
Format & System & Typical field & \code{style} value \\
\midrule
APA & Author-date & Psychology, education, social science & \code{apa} \\
Harvard & Author-date & Business, social science & \code{harvard-cite-them-right} \\
Chicago & Author-date & History, humanities & \code{chicago-author-date} \\
MLA & Author-page & Literature, languages, arts & \code{mla} \\
IEEE & Numeric & Computer science, electrical engineering & \code{ieee} \\
Vancouver & Numeric & Medicine, nursing & \code{vancouver} \\
\end{tlt}

Just change \code{style} to the value in the last column of the table, for example \code{"apa"} to \code{"ieee"}, and after recompiling the citations in the text and the list at the end switch to the matching style together, with no need to change how \code{\citep}, \code{\citet} and the others are written. After changing it, look at both the body and the reference page to check how the authors, years or numbers are shown.

The list above shows only the most common ones, and many more styles are available; if the format your department names is not among them, obtain the matching \code{.csl} style file (Bluebook, commonly used in law, is one example) and give its path in \code{style}, for example \code{style: "/misc/my-style.csl"}.

The default \code{full: false} lists only the works cited in the text. If a new entry still does not appear in the list, first check that the text cites it and that the capitalisation of its key matches. To list every entry of the bibliography file for the time being, change it to \code{true} and check the data entry by entry.

The place where the references are inserted is decided by \code{\references} in \code{main.tex}. Below it sits after Chapter 4 and before the appendices; if your department wants a different binding order, just move this line, with no need to copy the reference list by hand:

\begin{lstlisting}[language=latex]
\chapter{文獻、程式碼與附錄}\label{cha:refs}

本章接續說明文獻引用、程式碼與附錄的寫法。各節範例互不依賴，需要哪個就照抄哪段，再換成自己論文的資料即可。

文獻與附錄的檔案放在 \code{contents/back/}，其餘範例可以直接寫進本章。搬動或修改內容時，記得同步更新 \code{main.tex} 裡對應的引用位置，以免參考文獻或附錄的順序跑掉。

\section{引用文獻}

\code{.bib} 是以 \BibTeX{} 格式儲存的純文字文獻資料庫，一筆文獻對應一個條目，格式是 \code{@類型{識別名稱, 欄位 = {值}, ...}}。\code{contents/back/references.bib} 已經收錄幾筆範例，其中一筆長這樣：

\begin{lstlisting}
@inproceedings{kocher99,
  author =       {C. Kocher and J. Jaffe and B. Jun},
  title =        {Differential Power Analysis},
  editor =       {M. Wiener},
  booktitle =    {Advances in Cryptology ({CRYPTO}~'99)},
  series =       {Lecture Notes in Computer Science},
  volume =       1666,
  pages =        {388-397},
  publisher =    {Springer-Verlag},
  month =        {August},
  year =         1999
}
\end{lstlisting}

\code{@inproceedings} 是文獻類型（研討會論文），常見的還有 \code{@article}（期刊論文）、\code{@book}（書籍）、\code{@incollection}（專書篇章）等，類型會影響參考文獻頁該顯示哪些欄位。\code{kocher99} 是識別名稱，用來在內文引用這筆文獻，同一份 \code{.bib} 裡不能重複；後面的 \code{author}、\code{title}、\code{year} 等則是這筆文獻的資料，逐欄填入即可，順序不影響排版結果。

Google Scholar 每筆搜尋結果下方的「引用」按鈕，彈出視窗最下面就有 \BibTeX{} 連結，點開即可複製整個條目；圖書館資料庫（如 IEEE Xplore、ACM Digital Library）的文獻頁面也大多有「Export Citation」或「Cite」功能，選 \BibTeX{} 格式匯出。若平常使用 Zotero、EndNote 等文獻管理軟體，同樣可以直接匯出 \code{.bib} 檔。拿到條目後貼進 \code{references.bib}，把識別名稱改成自己好記的名字（例如 \code{作者姓+年份}），再確認沒有跟既有條目重複即可。

文獻資料準備好後，回到章節檔案，使用相同的識別名稱引用。若要把文獻放在句尾，可以使用 \code{\citep}；若要把作者寫進句子，則使用 \code{\citet}。以下也列出只顯示作者或年份的寫法，可以依句子需要選用。這四個指令來自 \code{citation-style-language} 套件，模板已經載入。

\begin{lstlisting}[language=latex]
\citep{kocher99}       % 括號引用
\citet{kocher99}       % 行文引用
\citeauthor{kocher99}  % 只要作者
\citeyear{kocher99}    % 只要年份
\end{lstlisting}

四個指令編譯後的差異如下，請先對照確認，再將識別名稱換成自己要引用的文獻：

\begin{tlt}{文獻引用指令對照}{tbl:cite-forms}{L{0.3\linewidth}L{0.25\linewidth}L{0.35\linewidth}}
指令 & 顯示結果 & 用途 \\
\midrule
\code{\citep{kocher99}} & \citep{kocher99} & 括號引用，放在句尾 \\
\code{\citet{kocher99}} & \citet{kocher99} & 行文引用，把作者寫進句子 \\
\code{\citeauthor{kocher99}} & \citeauthor{kocher99} & 只顯示作者 \\
\code{\citeyear{kocher99}} & \citeyear{kocher99} & 只顯示年份 \\
\end{tlt}

引用書中的特定段落時，可以用方括號補上頁碼或章節。以下以 \code{companion} 這筆文獻為例，加入「第 42 頁」；套用到自己的論文時，請改成實際引用的位置：

\begin{lstlisting}[language=latex]
\citep[第 42 頁]{companion}
\end{lstlisting}

編譯後會顯示為 \citep[第 42 頁]{companion}。頁碼會與文獻引用一起排版，標點和排列方式則由選用的參考文獻樣式決定。

若要一次引用多筆，把識別名稱用逗號隔開，例如 \code{\citep{kocher99,IEEE-1363}}。這種寫法與單筆引用一樣會使用設定檔中的文獻樣式，可以選擇自己較習慣的方式。

\section{參考文獻設定}

請打開 \code{config.yml}，在既有的 \code{bibliography:} 區塊下加入需要的設定。\code{path} 用來指定文獻檔案，\code{style} 決定引用格式；若資料分散在多個檔案，可以在 \code{path} 的清單中依序加入路徑。

\begin{lstlisting}[language=yaml]
bibliography:
  path: ["/contents/back/references.bib"] # 可以給多個檔案
  style: "apa"       # 內建 CSL 樣式名稱，或 .csl 檔案路徑，見下方對照表
  title: "參~考~文~獻"
  color: "#0000ff"   # 引用與連結的顏色，要黑白印就改成 "#000000"
  full: false        # true 會把沒引用到的條目也列出來
\end{lstlisting}

\code{path} 可以混用 \code{.bib}、Hayagriva \code{.yml}（Typst 版使用的格式）與 CSL JSON／YAML；Hayagriva 檔會在 \code{scripts/gen.py} 執行時自動轉成 CSL JSON，不必自己轉換。不同檔案裡的識別名稱不能重複。

\code{style} 決定文獻的排版邏輯，大致分成三個系統：作者-年份制在正文顯示作者姓名與年份，讀者能一眼看出研究新舊；作者-頁碼制在正文顯示作者與頁碼，方便對照原文出處；數字制則在正文只留數字或上標，清單依編號排列，讀起來最不受干擾。模板以 \code{citation-style-language} 與 \code{citeproc-lua} 處理 CSL 樣式，隨 \TeX{} 發行版安裝的樣式涵蓋各領域慣用的幾種，如下：

\begin{tlt}[notes={\texttt{mla} 是 \texttt{modern-language-association} 的簡稱，兩種寫法都可以。}]{常見文獻格式對照}{tbl:citation-styles}{L{0.12\linewidth}L{0.17\linewidth}L{0.25\linewidth}L{0.26\linewidth}}
格式 & 系統 & 常見領域 & \code{style} 值 \\
\midrule
APA & 作者-年份制 & 心理、教育、社會科學 & \code{apa} \\
Harvard & 作者-年份制 & 商管、社會科學 & \code{harvard-cite-them-right} \\
Chicago & 作者-年份制 & 歷史、人文 & \code{chicago-author-date} \\
MLA & 作者-頁碼制 & 文學、語言、藝術 & \code{mla} \\
IEEE & 數字制 & 資工、電機、工程 & \code{ieee} \\
Vancouver & 數字制 & 醫學、護理 & \code{vancouver} \\
\end{tlt}

\code{style} 直接改成對照表最後一欄的值即可，例如把 \code{"apa"} 改成 \code{"ieee"}，重新編譯後，正文引用與文末清單會一起換成對應樣式，就不用再額外調整 \code{\citep}、\code{\citet} 等指令的寫法。修改後請同時查看正文與參考文獻頁，確認作者、年份或編號的呈現方式。

以上清單只列出最常見的幾種，實際可用的樣式遠不止這些；若系所指定的格式不在其中（例如法律學門常用的 Bluebook），可以自行取得對應的 \code{.csl} 樣式檔，將 \code{style} 改填該檔案的路徑，例如 \code{style: "/misc/my-style.csl"}。

預設的 \code{full: false} 只會列出正文已引用的文獻。若加入新條目後，清單中仍看不到它，請先確認正文是否已經引用，以及識別名稱的大小寫是否一致。需要暫時列出文獻檔中的所有條目時，可以改為 \code{true}，方便逐筆檢查資料。

參考文獻的插入位置由 \code{main.tex} 中的 \code{\references} 決定。以下將它放在第四章之後、附錄之前；若系所要求不同的裝訂順序，移動此行即可，不需要手動複製文獻清單：

\begin{lstlisting}[language=latex]
\chapter{文獻、程式碼與附錄}\label{cha:refs}

本章接續說明文獻引用、程式碼與附錄的寫法。各節範例互不依賴，需要哪個就照抄哪段，再換成自己論文的資料即可。

文獻與附錄的檔案放在 \code{contents/back/}，其餘範例可以直接寫進本章。搬動或修改內容時，記得同步更新 \code{main.tex} 裡對應的引用位置，以免參考文獻或附錄的順序跑掉。

\section{引用文獻}

\code{.bib} 是以 \BibTeX{} 格式儲存的純文字文獻資料庫，一筆文獻對應一個條目，格式是 \code{@類型{識別名稱, 欄位 = {值}, ...}}。\code{contents/back/references.bib} 已經收錄幾筆範例，其中一筆長這樣：

\begin{lstlisting}
@inproceedings{kocher99,
  author =       {C. Kocher and J. Jaffe and B. Jun},
  title =        {Differential Power Analysis},
  editor =       {M. Wiener},
  booktitle =    {Advances in Cryptology ({CRYPTO}~'99)},
  series =       {Lecture Notes in Computer Science},
  volume =       1666,
  pages =        {388-397},
  publisher =    {Springer-Verlag},
  month =        {August},
  year =         1999
}
\end{lstlisting}

\code{@inproceedings} 是文獻類型（研討會論文），常見的還有 \code{@article}（期刊論文）、\code{@book}（書籍）、\code{@incollection}（專書篇章）等，類型會影響參考文獻頁該顯示哪些欄位。\code{kocher99} 是識別名稱，用來在內文引用這筆文獻，同一份 \code{.bib} 裡不能重複；後面的 \code{author}、\code{title}、\code{year} 等則是這筆文獻的資料，逐欄填入即可，順序不影響排版結果。

Google Scholar 每筆搜尋結果下方的「引用」按鈕，彈出視窗最下面就有 \BibTeX{} 連結，點開即可複製整個條目；圖書館資料庫（如 IEEE Xplore、ACM Digital Library）的文獻頁面也大多有「Export Citation」或「Cite」功能，選 \BibTeX{} 格式匯出。若平常使用 Zotero、EndNote 等文獻管理軟體，同樣可以直接匯出 \code{.bib} 檔。拿到條目後貼進 \code{references.bib}，把識別名稱改成自己好記的名字（例如 \code{作者姓+年份}），再確認沒有跟既有條目重複即可。

文獻資料準備好後，回到章節檔案，使用相同的識別名稱引用。若要把文獻放在句尾，可以使用 \code{\citep}；若要把作者寫進句子，則使用 \code{\citet}。以下也列出只顯示作者或年份的寫法，可以依句子需要選用。這四個指令來自 \code{citation-style-language} 套件，模板已經載入。

\begin{lstlisting}[language=latex]
\citep{kocher99}       % 括號引用
\citet{kocher99}       % 行文引用
\citeauthor{kocher99}  % 只要作者
\citeyear{kocher99}    % 只要年份
\end{lstlisting}

四個指令編譯後的差異如下，請先對照確認，再將識別名稱換成自己要引用的文獻：

\begin{tlt}{文獻引用指令對照}{tbl:cite-forms}{L{0.3\linewidth}L{0.25\linewidth}L{0.35\linewidth}}
指令 & 顯示結果 & 用途 \\
\midrule
\code{\citep{kocher99}} & \citep{kocher99} & 括號引用，放在句尾 \\
\code{\citet{kocher99}} & \citet{kocher99} & 行文引用，把作者寫進句子 \\
\code{\citeauthor{kocher99}} & \citeauthor{kocher99} & 只顯示作者 \\
\code{\citeyear{kocher99}} & \citeyear{kocher99} & 只顯示年份 \\
\end{tlt}

引用書中的特定段落時，可以用方括號補上頁碼或章節。以下以 \code{companion} 這筆文獻為例，加入「第 42 頁」；套用到自己的論文時，請改成實際引用的位置：

\begin{lstlisting}[language=latex]
\citep[第 42 頁]{companion}
\end{lstlisting}

編譯後會顯示為 \citep[第 42 頁]{companion}。頁碼會與文獻引用一起排版，標點和排列方式則由選用的參考文獻樣式決定。

若要一次引用多筆，把識別名稱用逗號隔開，例如 \code{\citep{kocher99,IEEE-1363}}。這種寫法與單筆引用一樣會使用設定檔中的文獻樣式，可以選擇自己較習慣的方式。

\section{參考文獻設定}

請打開 \code{config.yml}，在既有的 \code{bibliography:} 區塊下加入需要的設定。\code{path} 用來指定文獻檔案，\code{style} 決定引用格式；若資料分散在多個檔案，可以在 \code{path} 的清單中依序加入路徑。

\begin{lstlisting}[language=yaml]
bibliography:
  path: ["/contents/back/references.bib"] # 可以給多個檔案
  style: "apa"       # 內建 CSL 樣式名稱，或 .csl 檔案路徑，見下方對照表
  title: "參~考~文~獻"
  color: "#0000ff"   # 引用與連結的顏色，要黑白印就改成 "#000000"
  full: false        # true 會把沒引用到的條目也列出來
\end{lstlisting}

\code{path} 可以混用 \code{.bib}、Hayagriva \code{.yml}（Typst 版使用的格式）與 CSL JSON／YAML；Hayagriva 檔會在 \code{scripts/gen.py} 執行時自動轉成 CSL JSON，不必自己轉換。不同檔案裡的識別名稱不能重複。

\code{style} 決定文獻的排版邏輯，大致分成三個系統：作者-年份制在正文顯示作者姓名與年份，讀者能一眼看出研究新舊；作者-頁碼制在正文顯示作者與頁碼，方便對照原文出處；數字制則在正文只留數字或上標，清單依編號排列，讀起來最不受干擾。模板以 \code{citation-style-language} 與 \code{citeproc-lua} 處理 CSL 樣式，隨 \TeX{} 發行版安裝的樣式涵蓋各領域慣用的幾種，如下：

\begin{tlt}[notes={\texttt{mla} 是 \texttt{modern-language-association} 的簡稱，兩種寫法都可以。}]{常見文獻格式對照}{tbl:citation-styles}{L{0.12\linewidth}L{0.17\linewidth}L{0.25\linewidth}L{0.26\linewidth}}
格式 & 系統 & 常見領域 & \code{style} 值 \\
\midrule
APA & 作者-年份制 & 心理、教育、社會科學 & \code{apa} \\
Harvard & 作者-年份制 & 商管、社會科學 & \code{harvard-cite-them-right} \\
Chicago & 作者-年份制 & 歷史、人文 & \code{chicago-author-date} \\
MLA & 作者-頁碼制 & 文學、語言、藝術 & \code{mla} \\
IEEE & 數字制 & 資工、電機、工程 & \code{ieee} \\
Vancouver & 數字制 & 醫學、護理 & \code{vancouver} \\
\end{tlt}

\code{style} 直接改成對照表最後一欄的值即可，例如把 \code{"apa"} 改成 \code{"ieee"}，重新編譯後，正文引用與文末清單會一起換成對應樣式，就不用再額外調整 \code{\citep}、\code{\citet} 等指令的寫法。修改後請同時查看正文與參考文獻頁，確認作者、年份或編號的呈現方式。

以上清單只列出最常見的幾種，實際可用的樣式遠不止這些；若系所指定的格式不在其中（例如法律學門常用的 Bluebook），可以自行取得對應的 \code{.csl} 樣式檔，將 \code{style} 改填該檔案的路徑，例如 \code{style: "/misc/my-style.csl"}。

預設的 \code{full: false} 只會列出正文已引用的文獻。若加入新條目後，清單中仍看不到它，請先確認正文是否已經引用，以及識別名稱的大小寫是否一致。需要暫時列出文獻檔中的所有條目時，可以改為 \code{true}，方便逐筆檢查資料。

參考文獻的插入位置由 \code{main.tex} 中的 \code{\references} 決定。以下將它放在第四章之後、附錄之前；若系所要求不同的裝訂順序，移動此行即可，不需要手動複製文獻清單：

\begin{lstlisting}[language=latex]
\chapter{文獻、程式碼與附錄}\label{cha:refs}

本章接續說明文獻引用、程式碼與附錄的寫法。各節範例互不依賴，需要哪個就照抄哪段，再換成自己論文的資料即可。

文獻與附錄的檔案放在 \code{contents/back/}，其餘範例可以直接寫進本章。搬動或修改內容時，記得同步更新 \code{main.tex} 裡對應的引用位置，以免參考文獻或附錄的順序跑掉。

\section{引用文獻}

\code{.bib} 是以 \BibTeX{} 格式儲存的純文字文獻資料庫，一筆文獻對應一個條目，格式是 \code{@類型{識別名稱, 欄位 = {值}, ...}}。\code{contents/back/references.bib} 已經收錄幾筆範例，其中一筆長這樣：

\begin{lstlisting}
@inproceedings{kocher99,
  author =       {C. Kocher and J. Jaffe and B. Jun},
  title =        {Differential Power Analysis},
  editor =       {M. Wiener},
  booktitle =    {Advances in Cryptology ({CRYPTO}~'99)},
  series =       {Lecture Notes in Computer Science},
  volume =       1666,
  pages =        {388-397},
  publisher =    {Springer-Verlag},
  month =        {August},
  year =         1999
}
\end{lstlisting}

\code{@inproceedings} 是文獻類型（研討會論文），常見的還有 \code{@article}（期刊論文）、\code{@book}（書籍）、\code{@incollection}（專書篇章）等，類型會影響參考文獻頁該顯示哪些欄位。\code{kocher99} 是識別名稱，用來在內文引用這筆文獻，同一份 \code{.bib} 裡不能重複；後面的 \code{author}、\code{title}、\code{year} 等則是這筆文獻的資料，逐欄填入即可，順序不影響排版結果。

Google Scholar 每筆搜尋結果下方的「引用」按鈕，彈出視窗最下面就有 \BibTeX{} 連結，點開即可複製整個條目；圖書館資料庫（如 IEEE Xplore、ACM Digital Library）的文獻頁面也大多有「Export Citation」或「Cite」功能，選 \BibTeX{} 格式匯出。若平常使用 Zotero、EndNote 等文獻管理軟體，同樣可以直接匯出 \code{.bib} 檔。拿到條目後貼進 \code{references.bib}，把識別名稱改成自己好記的名字（例如 \code{作者姓+年份}），再確認沒有跟既有條目重複即可。

文獻資料準備好後，回到章節檔案，使用相同的識別名稱引用。若要把文獻放在句尾，可以使用 \code{\citep}；若要把作者寫進句子，則使用 \code{\citet}。以下也列出只顯示作者或年份的寫法，可以依句子需要選用。這四個指令來自 \code{citation-style-language} 套件，模板已經載入。

\begin{lstlisting}[language=latex]
\citep{kocher99}       % 括號引用
\citet{kocher99}       % 行文引用
\citeauthor{kocher99}  % 只要作者
\citeyear{kocher99}    % 只要年份
\end{lstlisting}

四個指令編譯後的差異如下，請先對照確認，再將識別名稱換成自己要引用的文獻：

\begin{tlt}{文獻引用指令對照}{tbl:cite-forms}{L{0.3\linewidth}L{0.25\linewidth}L{0.35\linewidth}}
指令 & 顯示結果 & 用途 \\
\midrule
\code{\citep{kocher99}} & \citep{kocher99} & 括號引用，放在句尾 \\
\code{\citet{kocher99}} & \citet{kocher99} & 行文引用，把作者寫進句子 \\
\code{\citeauthor{kocher99}} & \citeauthor{kocher99} & 只顯示作者 \\
\code{\citeyear{kocher99}} & \citeyear{kocher99} & 只顯示年份 \\
\end{tlt}

引用書中的特定段落時，可以用方括號補上頁碼或章節。以下以 \code{companion} 這筆文獻為例，加入「第 42 頁」；套用到自己的論文時，請改成實際引用的位置：

\begin{lstlisting}[language=latex]
\citep[第 42 頁]{companion}
\end{lstlisting}

編譯後會顯示為 \citep[第 42 頁]{companion}。頁碼會與文獻引用一起排版，標點和排列方式則由選用的參考文獻樣式決定。

若要一次引用多筆，把識別名稱用逗號隔開，例如 \code{\citep{kocher99,IEEE-1363}}。這種寫法與單筆引用一樣會使用設定檔中的文獻樣式，可以選擇自己較習慣的方式。

\section{參考文獻設定}

請打開 \code{config.yml}，在既有的 \code{bibliography:} 區塊下加入需要的設定。\code{path} 用來指定文獻檔案，\code{style} 決定引用格式；若資料分散在多個檔案，可以在 \code{path} 的清單中依序加入路徑。

\begin{lstlisting}[language=yaml]
bibliography:
  path: ["/contents/back/references.bib"] # 可以給多個檔案
  style: "apa"       # 內建 CSL 樣式名稱，或 .csl 檔案路徑，見下方對照表
  title: "參~考~文~獻"
  color: "#0000ff"   # 引用與連結的顏色，要黑白印就改成 "#000000"
  full: false        # true 會把沒引用到的條目也列出來
\end{lstlisting}

\code{path} 可以混用 \code{.bib}、Hayagriva \code{.yml}（Typst 版使用的格式）與 CSL JSON／YAML；Hayagriva 檔會在 \code{scripts/gen.py} 執行時自動轉成 CSL JSON，不必自己轉換。不同檔案裡的識別名稱不能重複。

\code{style} 決定文獻的排版邏輯，大致分成三個系統：作者-年份制在正文顯示作者姓名與年份，讀者能一眼看出研究新舊；作者-頁碼制在正文顯示作者與頁碼，方便對照原文出處；數字制則在正文只留數字或上標，清單依編號排列，讀起來最不受干擾。模板以 \code{citation-style-language} 與 \code{citeproc-lua} 處理 CSL 樣式，隨 \TeX{} 發行版安裝的樣式涵蓋各領域慣用的幾種，如下：

\begin{tlt}[notes={\texttt{mla} 是 \texttt{modern-language-association} 的簡稱，兩種寫法都可以。}]{常見文獻格式對照}{tbl:citation-styles}{L{0.12\linewidth}L{0.17\linewidth}L{0.25\linewidth}L{0.26\linewidth}}
格式 & 系統 & 常見領域 & \code{style} 值 \\
\midrule
APA & 作者-年份制 & 心理、教育、社會科學 & \code{apa} \\
Harvard & 作者-年份制 & 商管、社會科學 & \code{harvard-cite-them-right} \\
Chicago & 作者-年份制 & 歷史、人文 & \code{chicago-author-date} \\
MLA & 作者-頁碼制 & 文學、語言、藝術 & \code{mla} \\
IEEE & 數字制 & 資工、電機、工程 & \code{ieee} \\
Vancouver & 數字制 & 醫學、護理 & \code{vancouver} \\
\end{tlt}

\code{style} 直接改成對照表最後一欄的值即可，例如把 \code{"apa"} 改成 \code{"ieee"}，重新編譯後，正文引用與文末清單會一起換成對應樣式，就不用再額外調整 \code{\citep}、\code{\citet} 等指令的寫法。修改後請同時查看正文與參考文獻頁，確認作者、年份或編號的呈現方式。

以上清單只列出最常見的幾種，實際可用的樣式遠不止這些；若系所指定的格式不在其中（例如法律學門常用的 Bluebook），可以自行取得對應的 \code{.csl} 樣式檔，將 \code{style} 改填該檔案的路徑，例如 \code{style: "/misc/my-style.csl"}。

預設的 \code{full: false} 只會列出正文已引用的文獻。若加入新條目後，清單中仍看不到它，請先確認正文是否已經引用，以及識別名稱的大小寫是否一致。需要暫時列出文獻檔中的所有條目時，可以改為 \code{true}，方便逐筆檢查資料。

參考文獻的插入位置由 \code{main.tex} 中的 \code{\references} 決定。以下將它放在第四章之後、附錄之前；若系所要求不同的裝訂順序，移動此行即可，不需要手動複製文獻清單：

\begin{lstlisting}[language=latex]
\input{contents/chapter04}

\references                          % 參考文獻

\appendix
\input{contents/back/appendix01}     % 附錄
\end{lstlisting}

\section{程式碼區塊}

要在論文中列出程式碼，可以直接在章節檔案裡建立程式碼區塊。模板會套用等寬字型與淺灰底色，並依指定的語言進行語法標色。先選擇符合內容的語言名稱，例如 \code{python}、\code{bash} 或 \code{yaml} 等，再填入程式碼即可。

建立程式碼區塊的方法是用 \code{lstlisting} 環境包住需要顯示的程式碼，並在方括號裡用 \code{language=} 指定語言名稱。以 Python 程式的寫法為例：

\begin{codeblock}[language=latex]
\begin{lstlisting}[language=python]
def solve(n):
    return sum(range(n))
\end{lstlisting}
\end{codeblock}

如果這段程式需要標題與編號，在方括號裡加上 \code{caption=} 與 \code{label=}。需要從內文引用時，用 \code{\lstref} 引用，寫法和前一章的圖表相同：

\begin{codeblock}[language=latex]
\begin{lstlisting}[language=python,caption={遞迴版本},label=lst:fib]
def fib(n):
    return n if n < 2 else fib(n - 1) + fib(n - 2)
\end{lstlisting}
\end{codeblock}

注意到：程式碼預設不顯示標題與編號，下面也採用不加標題的寫法。需要標題時，可以依上面的範例選填 \code{caption}，並用 \code{\lstref{lst:fib}} 引用加上標籤的程式碼。標題預設放在下方；若要移到上方，請將 \code{layout.caption-position.code} 設為 \code{"top"}。

\begin{lstlisting}[language=python]
def fib(n):
    return n if n < 2 else fib(n - 1) + fib(n - 2)
\end{lstlisting}

當程式碼超過一頁時，會與表格一樣自動跨頁，不必手動拆成數個區塊——直接使用 \code{lstlisting} 的區塊，以及加上標題的區塊，都支援跨頁；有標題時，整段程式仍共用一個編號與標題。若要把外部檔案的內容直接放進論文，可以用 \code{\lstinputlisting[language=python]{path/to/file.py}}。

若只是在句子中提到指令或檔名，請用 \code{\code} 包住文字，例如

\begin{lstlisting}[language=latex]
\code{make release}
\end{lstlisting}

這種行內寫法會保留等寬字型，字級與基線則和本文一致，不會另外加上底色或框線。因為它是逐字排版，不能寫在 \code{\section} 或 \code{\caption} 這類的大括號引數裡；那些地方請改用 \code{\texttt}，需要顯示反斜線等特殊符號時，再加上 \code{\textbackslash}。

如果需要調整程式碼區塊外觀時，請打開 \code{tw-thesis.sty}，修改程式碼區塊那一節最上方的主題設定即可。此處的設定會套用到共用的程式碼樣式，不必逐一修改各章範例；其中 \code{\twcodesize} 控制區塊字級，行內程式碼仍會跟隨所在段落的字級。以下列出各項設定：

\begin{lstlisting}[language=latex]
\definecolor{twcodebg}{gray}{0.973}      % 底色
\definecolor{twcodeblue}{HTML}{4B69C6}   % 行內程式碼與關鍵字的顏色
\newcommand\twcodesize{\fontsize{11.4}{15}\selectfont}   % 區塊字級
\setmonofont{DejaVu Sans Mono}           % 等寬字體，可以放自己的字型
\end{lstlisting}

\begin{tlt}{程式碼樣式參數}{tbl:code-theme}{L{0.28\linewidth}L{0.62\linewidth}}
參數 & 用途與範例 \\
\midrule
\code{\setmonofont} & 等寬字型；例：\code{\setmonofont{DejaVu Sans Mono}}。 \\
\code{\twcodesize} & 區塊字級；例：本文的 0.95 倍，行內程式碼仍跟隨本文。 \\
\code{twcodebg} & 區塊底色；例：\code{\definecolor{twcodebg}{gray}{0.973}}。 \\
\code{twcodeblue} & 行內程式碼與關鍵字的顏色；例：\code{4B69C6}。 \\
\code{twcodegray} & 註解的顏色。 \\
\end{tlt}

\section{附錄}

需要收錄問卷、補充資料或較長的推導時，可以放到 \code{contents/back/} 的附錄檔案。請先在 \code{main.tex} 寫一次 \code{\appendix}，再填入標題與內文。以下示範第一份附錄的寫法，編譯後會使用 A、B 等字母編號：

\begin{lstlisting}[language=latex]
\chapter{問卷內容}\label{appendix-a}

\section{第一部分}
\end{lstlisting}

\code{main.tex} 裡的 \code{\appendix} 會切換編號方式，並把章的計數歸零。因此接下來的章標題會顯示為「附錄 A」，節標題則使用 \code{A.1}。替換內容時，只需修改「問卷內容」與各節標題，不必在標題中手動輸入字母或編號。

若要新增第二份附錄，請建立另一個 \code{.tex} 檔案，同樣用 \code{\chapter} 寫標題，但不要再寫一次 \code{\appendix}。接著到 \code{main.tex}，在第一份附錄後面加入新檔案的 \code{\input}，編號就會接續為 B。

\code{\appendix} 之後編號的規則是固定的，第二份、第三份附錄仍使用相同設定，不要自己改成 \code{B.1} 或 \code{C.1}。

\section{學校規範更新}

各校的論文格式規範偶爾會修訂，而模板裡的 \code{schools/*.yml} 和 \code{docs/} 收錄的規範 PDF，都只是某個時間點的版本。為了避免在不知情的情況下照著舊規範排版，搭配模板使用的管理工具 \code{tw-thesis} 會定時抓取是否有新的論文規範。

\subsection{提醒運作方式}

另一個專案 \href{https://github.com/TW-Thesis/tw-thesis-updater}{\texttt{tw-thesis-updater}} 每週自動抓取各校的官方規範文件，並公開一份清單，記錄每間學校目前最新版本的日期與雜湊值。在專案資料夾內執行 \code{tw-thesis} 開啟管理介面後，首頁會拿 \code{config.yml} 選定的學校，比對專案 \code{docs/sources.lock.json} 裡記錄的版本；兩者不同便會出現黃色提醒框，說明該校規範已修訂、附上新規範的下載連結。

\begin{lstlisting}[language=bash]
tw-thesis    # 在專案資料夾內執行；首頁會顯示提醒
\end{lstlisting}

該檢查只針對選擇的那所學校，一天最多連線一次；離線、或清單尚未發布時不會顯示任何錯誤，也不會影響編譯。

\subsection{提醒後續事宜}

程式僅負責提醒，不會自動修改任何排版設定，也不會判斷新規範和舊規範差在哪裡：各校、各系所對規範的解讀不盡相同，有些條文（例如封面格式、裝訂順序）尚要符合系所額外的規定，所以是否要修改、修改幅度仍由使用者決定。看到提醒後，建議依序處理：

\begin{enumerate}
\item 從提醒框的連結下載新規範，並實際從頭讀過一遍，切勿只看公告的標題。如果學校網站同時提供「修正對照表」或「修訂說明」，先看那份就能最快能找出有哪幾條規則修改了。
\item 對照 \code{schools/<學校代碼>.yml}（例如政大是 \code{schools/nccu.yml}），逐項確認邊界、字級、行距、標題格式、封面與審定書、頁碼、參考文獻格式等設定是否需要調整。各項目的設定方式見\ref{cha:layout}；若規範要求的項目在學校設定檔裡找不到對應欄位，先到 \code{schools/default.yml} 查有沒有同名的設定。
\item 想先試試看、不想直接動學校設定檔時，可以把要調整的項目寫進自己的 \code{config.yml}。它的優先順序高於學校設定，確認結果沒問題後，再決定要不要正式改 \code{schools/}。
\item 需要的話，把新的 PDF 放進 \code{docs/} 取代舊檔（檔名內含版本日期），以後想對照時才不會拿到舊版本。
\item 調整完成後重新編譯，逐頁核對封面、審定書、目次、正文第一頁與參考文獻頁，最好再請指導教授或系辦幫忙確認一次。
\end{enumerate}

\subsection{免責聲明}

本模板的格式設定僅供參考，實際仍請以學校公布的論文規範及系所要求為準。規定往往不只一份：學校公布的是通則，系所、學院或指導教授另有補充要求（例如封面、裝訂、字體或圖表格式）的情況並不少見，這些都不會出現在模板的設定裡。

此外，規範更新提醒是靠定期抓取學校網站，可能晚於學校實際公告才出現，也可能因為學校改版網站而暫時抓不到，所以沒有收到提醒，並不代表規範沒有修訂。提交論文前，請務必再向系所或圖書館確認最新要求。


\references                          % 參考文獻

\appendix
\chapter{編譯問題排查}\label{appendix-a}

使用模板時，若遇到編譯失敗、圖片找不到或 PDF 沒有更新，可以先依本附錄的順序檢查。建議保留終端機中的錯誤訊息，從第一個錯誤指出的檔案與位置開始處理，修正後再重新編譯。

這份檔案也示範附錄的排版方式。換成自己的論文時，可以保留章標題的寫法，將標題與內文替換成問卷、補充數據或其他資料，再到 \code{main.tex} 確認附錄的裝訂順序。

\section{指令執行失敗}

若終端機顯示找不到 \code{xelatex} 或 \code{latexmk}，請先確認已完成安裝，再重新開啟終端機，執行

\begin{lstlisting}[language=bash]
xelatex --version
\end{lstlisting}

確認能夠顯示版本資訊後，再嘗試編譯。若找不到的是 \code{make}，則可以先使用以下指令，直接產生 PDF：

\begin{lstlisting}[language=bash]
texlua scripts/gen.lua
latexmk -xelatex main.tex
\end{lstlisting}

若錯誤訊息指出找不到 \code{main.tex} 或 \code{.tw-thesis-config.tex}，請檢查終端機目前所在的資料夾。編譯指令應在專案根目錄執行，也就是能看到 \code{main.tex}、\code{config.yml} 與 \code{contents/} 的位置；\code{.tw-thesis-config.tex} 則是 \code{scripts/gen.lua} 產生的，先執行過一次就會出現。若目前位於章節資料夾，請先切換回根目錄再試一次。

\section{修改設定後無法編譯}

修改設定後出現錯誤時，請先看訊息指出哪個檔案、哪一個欄位。若包含 \code{unknown key}，可以對照 \code{schools/default.yml} 檢查欄位拼字與所在層級；若是 \code{YAML} 解析錯誤，則先確認縮排使用空白，且冒號後保留空白，再檢查引號是否成對。

加入新設定時，請先找出檔案中是否已有相同區塊，避免在同一層重複建立 \code{layout:}。例如要同時修改字級與左側邊界，應把兩項設定放在同一個 \code{layout:} 下，並讓 \code{left} 再往內縮排一層：

\begin{lstlisting}[language=yaml]
layout:
  font-size: 12
  margin:
    left: 3.5
\end{lstlisting}

若一次修改了多個設定，暫時無法判斷錯誤來源，可以先將剛加入的內容複製到別處保存，再移除這些設定，確認原本的文件能正常編譯。之後逐項加回，每次存檔後檢查結果，較容易找到需要修正的位置。

\section{圖片或引用找不到}

圖片無法載入時，請先確認檔案確實放在 \code{figures/} 的指定位置，再比對路徑中的檔名、副檔名與大小寫。路徑是從專案根目錄算起的，不是從目前章節所在的資料夾開始。若曾經移動或重新命名圖片，也要一併修改章節中的路徑。

圖表引用顯示成 \code{??} 時，請檢查圖表是否已加上 \code{\label}，以及名稱是否與內文的 \code{\ref} 完全相同；也別忘了 \LaTeX{} 需要編譯兩次才會解出交叉引用，\code{latexmk} 會自動處理。若是文獻引用出錯，則到 \code{.bib} 檔案確認識別名稱，再檢查 \code{config.yml} 的 \code{bibliography.path} 是否包含該文獻檔。修正後重新編譯，確認引用已顯示為編號或文獻資料。

\section{PDF 沒有更新}

修改內容後，請先確認終端機是否已完成這次編譯，以及是否出現錯誤。編譯失敗時，資料夾中可能仍保留上一次的 \code{main.pdf}，因此能開啟 PDF 並不代表新內容已經寫入。請先修正錯誤，再重新執行編譯指令。

若使用 \code{make watch}，請確認監看程序仍在執行，並且修改的檔案已經存檔。編譯成功後畫面仍未更新時，可以關閉再重新開啟 PDF，同時確認檢視器開啟的是目前專案中的 \code{main.pdf}，避免看到先前複製到其他位置的版本。若編譯結果卡在舊的狀態，執行 \code{make clean} 清掉暫存檔後再編譯一次，通常就能解決。
     % 附錄
\end{lstlisting}

\section{程式碼區塊}

要在論文中列出程式碼，可以直接在章節檔案裡建立程式碼區塊。模板會套用等寬字型與淺灰底色，並依指定的語言進行語法標色。先選擇符合內容的語言名稱，例如 \code{python}、\code{bash} 或 \code{yaml} 等，再填入程式碼即可。

建立程式碼區塊的方法是用 \code{lstlisting} 環境包住需要顯示的程式碼，並在方括號裡用 \code{language=} 指定語言名稱。以 Python 程式的寫法為例：

\begin{codeblock}[language=latex]
\begin{lstlisting}[language=python]
def solve(n):
    return sum(range(n))
\end{lstlisting}
\end{codeblock}

如果這段程式需要標題與編號，在方括號裡加上 \code{caption=} 與 \code{label=}。需要從內文引用時，用 \code{\lstref} 引用，寫法和前一章的圖表相同：

\begin{codeblock}[language=latex]
\begin{lstlisting}[language=python,caption={遞迴版本},label=lst:fib]
def fib(n):
    return n if n < 2 else fib(n - 1) + fib(n - 2)
\end{lstlisting}
\end{codeblock}

注意到：程式碼預設不顯示標題與編號，下面也採用不加標題的寫法。需要標題時，可以依上面的範例選填 \code{caption}，並用 \code{\lstref{lst:fib}} 引用加上標籤的程式碼。標題預設放在下方；若要移到上方，請將 \code{layout.caption-position.code} 設為 \code{"top"}。

\begin{lstlisting}[language=python]
def fib(n):
    return n if n < 2 else fib(n - 1) + fib(n - 2)
\end{lstlisting}

當程式碼超過一頁時，會與表格一樣自動跨頁，不必手動拆成數個區塊——直接使用 \code{lstlisting} 的區塊，以及加上標題的區塊，都支援跨頁；有標題時，整段程式仍共用一個編號與標題。若要把外部檔案的內容直接放進論文，可以用 \code{\lstinputlisting[language=python]{path/to/file.py}}。

若只是在句子中提到指令或檔名，請用 \code{\code} 包住文字，例如

\begin{lstlisting}[language=latex]
\code{make release}
\end{lstlisting}

這種行內寫法會保留等寬字型，字級與基線則和本文一致，不會另外加上底色或框線。因為它是逐字排版，不能寫在 \code{\section} 或 \code{\caption} 這類的大括號引數裡；那些地方請改用 \code{\texttt}，需要顯示反斜線等特殊符號時，再加上 \code{\textbackslash}。

如果需要調整程式碼區塊外觀時，請打開 \code{tw-thesis.sty}，修改程式碼區塊那一節最上方的主題設定即可。此處的設定會套用到共用的程式碼樣式，不必逐一修改各章範例；其中 \code{\twcodesize} 控制區塊字級，行內程式碼仍會跟隨所在段落的字級。以下列出各項設定：

\begin{lstlisting}[language=latex]
\definecolor{twcodebg}{gray}{0.973}      % 底色
\definecolor{twcodeblue}{HTML}{4B69C6}   % 行內程式碼與關鍵字的顏色
\newcommand\twcodesize{\fontsize{11.4}{15}\selectfont}   % 區塊字級
\setmonofont{DejaVu Sans Mono}           % 等寬字體，可以放自己的字型
\end{lstlisting}

\begin{tlt}{程式碼樣式參數}{tbl:code-theme}{L{0.28\linewidth}L{0.62\linewidth}}
參數 & 用途與範例 \\
\midrule
\code{\setmonofont} & 等寬字型；例：\code{\setmonofont{DejaVu Sans Mono}}。 \\
\code{\twcodesize} & 區塊字級；例：本文的 0.95 倍，行內程式碼仍跟隨本文。 \\
\code{twcodebg} & 區塊底色；例：\code{\definecolor{twcodebg}{gray}{0.973}}。 \\
\code{twcodeblue} & 行內程式碼與關鍵字的顏色；例：\code{4B69C6}。 \\
\code{twcodegray} & 註解的顏色。 \\
\end{tlt}

\section{附錄}

需要收錄問卷、補充資料或較長的推導時，可以放到 \code{contents/back/} 的附錄檔案。請先在 \code{main.tex} 寫一次 \code{\appendix}，再填入標題與內文。以下示範第一份附錄的寫法，編譯後會使用 A、B 等字母編號：

\begin{lstlisting}[language=latex]
\chapter{問卷內容}\label{appendix-a}

\section{第一部分}
\end{lstlisting}

\code{main.tex} 裡的 \code{\appendix} 會切換編號方式，並把章的計數歸零。因此接下來的章標題會顯示為「附錄 A」，節標題則使用 \code{A.1}。替換內容時，只需修改「問卷內容」與各節標題，不必在標題中手動輸入字母或編號。

若要新增第二份附錄，請建立另一個 \code{.tex} 檔案，同樣用 \code{\chapter} 寫標題，但不要再寫一次 \code{\appendix}。接著到 \code{main.tex}，在第一份附錄後面加入新檔案的 \code{\input}，編號就會接續為 B。

\code{\appendix} 之後編號的規則是固定的，第二份、第三份附錄仍使用相同設定，不要自己改成 \code{B.1} 或 \code{C.1}。

\section{學校規範更新}

各校的論文格式規範偶爾會修訂，而模板裡的 \code{schools/*.yml} 和 \code{docs/} 收錄的規範 PDF，都只是某個時間點的版本。為了避免在不知情的情況下照著舊規範排版，搭配模板使用的管理工具 \code{tw-thesis} 會定時抓取是否有新的論文規範。

\subsection{提醒運作方式}

另一個專案 \href{https://github.com/TW-Thesis/tw-thesis-updater}{\texttt{tw-thesis-updater}} 每週自動抓取各校的官方規範文件，並公開一份清單，記錄每間學校目前最新版本的日期與雜湊值。在專案資料夾內執行 \code{tw-thesis} 開啟管理介面後，首頁會拿 \code{config.yml} 選定的學校，比對專案 \code{docs/sources.lock.json} 裡記錄的版本；兩者不同便會出現黃色提醒框，說明該校規範已修訂、附上新規範的下載連結。

\begin{lstlisting}[language=bash]
tw-thesis    # 在專案資料夾內執行；首頁會顯示提醒
\end{lstlisting}

該檢查只針對選擇的那所學校，一天最多連線一次；離線、或清單尚未發布時不會顯示任何錯誤，也不會影響編譯。

\subsection{提醒後續事宜}

程式僅負責提醒，不會自動修改任何排版設定，也不會判斷新規範和舊規範差在哪裡：各校、各系所對規範的解讀不盡相同，有些條文（例如封面格式、裝訂順序）尚要符合系所額外的規定，所以是否要修改、修改幅度仍由使用者決定。看到提醒後，建議依序處理：

\begin{enumerate}
\item 從提醒框的連結下載新規範，並實際從頭讀過一遍，切勿只看公告的標題。如果學校網站同時提供「修正對照表」或「修訂說明」，先看那份就能最快能找出有哪幾條規則修改了。
\item 對照 \code{schools/<學校代碼>.yml}（例如政大是 \code{schools/nccu.yml}），逐項確認邊界、字級、行距、標題格式、封面與審定書、頁碼、參考文獻格式等設定是否需要調整。各項目的設定方式見\ref{cha:layout}；若規範要求的項目在學校設定檔裡找不到對應欄位，先到 \code{schools/default.yml} 查有沒有同名的設定。
\item 想先試試看、不想直接動學校設定檔時，可以把要調整的項目寫進自己的 \code{config.yml}。它的優先順序高於學校設定，確認結果沒問題後，再決定要不要正式改 \code{schools/}。
\item 需要的話，把新的 PDF 放進 \code{docs/} 取代舊檔（檔名內含版本日期），以後想對照時才不會拿到舊版本。
\item 調整完成後重新編譯，逐頁核對封面、審定書、目次、正文第一頁與參考文獻頁，最好再請指導教授或系辦幫忙確認一次。
\end{enumerate}

\subsection{免責聲明}

本模板的格式設定僅供參考，實際仍請以學校公布的論文規範及系所要求為準。規定往往不只一份：學校公布的是通則，系所、學院或指導教授另有補充要求（例如封面、裝訂、字體或圖表格式）的情況並不少見，這些都不會出現在模板的設定裡。

此外，規範更新提醒是靠定期抓取學校網站，可能晚於學校實際公告才出現，也可能因為學校改版網站而暫時抓不到，所以沒有收到提醒，並不代表規範沒有修訂。提交論文前，請務必再向系所或圖書館確認最新要求。


\references                          % 參考文獻

\appendix
\chapter{編譯問題排查}\label{appendix-a}

使用模板時，若遇到編譯失敗、圖片找不到或 PDF 沒有更新，可以先依本附錄的順序檢查。建議保留終端機中的錯誤訊息，從第一個錯誤指出的檔案與位置開始處理，修正後再重新編譯。

這份檔案也示範附錄的排版方式。換成自己的論文時，可以保留章標題的寫法，將標題與內文替換成問卷、補充數據或其他資料，再到 \code{main.tex} 確認附錄的裝訂順序。

\section{指令執行失敗}

若終端機顯示找不到 \code{xelatex} 或 \code{latexmk}，請先確認已完成安裝，再重新開啟終端機，執行

\begin{lstlisting}[language=bash]
xelatex --version
\end{lstlisting}

確認能夠顯示版本資訊後，再嘗試編譯。若找不到的是 \code{make}，則可以先使用以下指令，直接產生 PDF：

\begin{lstlisting}[language=bash]
texlua scripts/gen.lua
latexmk -xelatex main.tex
\end{lstlisting}

若錯誤訊息指出找不到 \code{main.tex} 或 \code{.tw-thesis-config.tex}，請檢查終端機目前所在的資料夾。編譯指令應在專案根目錄執行，也就是能看到 \code{main.tex}、\code{config.yml} 與 \code{contents/} 的位置；\code{.tw-thesis-config.tex} 則是 \code{scripts/gen.lua} 產生的，先執行過一次就會出現。若目前位於章節資料夾，請先切換回根目錄再試一次。

\section{修改設定後無法編譯}

修改設定後出現錯誤時，請先看訊息指出哪個檔案、哪一個欄位。若包含 \code{unknown key}，可以對照 \code{schools/default.yml} 檢查欄位拼字與所在層級；若是 \code{YAML} 解析錯誤，則先確認縮排使用空白，且冒號後保留空白，再檢查引號是否成對。

加入新設定時，請先找出檔案中是否已有相同區塊，避免在同一層重複建立 \code{layout:}。例如要同時修改字級與左側邊界，應把兩項設定放在同一個 \code{layout:} 下，並讓 \code{left} 再往內縮排一層：

\begin{lstlisting}[language=yaml]
layout:
  font-size: 12
  margin:
    left: 3.5
\end{lstlisting}

若一次修改了多個設定，暫時無法判斷錯誤來源，可以先將剛加入的內容複製到別處保存，再移除這些設定，確認原本的文件能正常編譯。之後逐項加回，每次存檔後檢查結果，較容易找到需要修正的位置。

\section{圖片或引用找不到}

圖片無法載入時，請先確認檔案確實放在 \code{figures/} 的指定位置，再比對路徑中的檔名、副檔名與大小寫。路徑是從專案根目錄算起的，不是從目前章節所在的資料夾開始。若曾經移動或重新命名圖片，也要一併修改章節中的路徑。

圖表引用顯示成 \code{??} 時，請檢查圖表是否已加上 \code{\label}，以及名稱是否與內文的 \code{\ref} 完全相同；也別忘了 \LaTeX{} 需要編譯兩次才會解出交叉引用，\code{latexmk} 會自動處理。若是文獻引用出錯，則到 \code{.bib} 檔案確認識別名稱，再檢查 \code{config.yml} 的 \code{bibliography.path} 是否包含該文獻檔。修正後重新編譯，確認引用已顯示為編號或文獻資料。

\section{PDF 沒有更新}

修改內容後，請先確認終端機是否已完成這次編譯，以及是否出現錯誤。編譯失敗時，資料夾中可能仍保留上一次的 \code{main.pdf}，因此能開啟 PDF 並不代表新內容已經寫入。請先修正錯誤，再重新執行編譯指令。

若使用 \code{make watch}，請確認監看程序仍在執行，並且修改的檔案已經存檔。編譯成功後畫面仍未更新時，可以關閉再重新開啟 PDF，同時確認檢視器開啟的是目前專案中的 \code{main.pdf}，避免看到先前複製到其他位置的版本。若編譯結果卡在舊的狀態，執行 \code{make clean} 清掉暫存檔後再編譯一次，通常就能解決。
     % 附錄
\end{lstlisting}

\section{程式碼區塊}

要在論文中列出程式碼，可以直接在章節檔案裡建立程式碼區塊。模板會套用等寬字型與淺灰底色，並依指定的語言進行語法標色。先選擇符合內容的語言名稱，例如 \code{python}、\code{bash} 或 \code{yaml} 等，再填入程式碼即可。

建立程式碼區塊的方法是用 \code{lstlisting} 環境包住需要顯示的程式碼，並在方括號裡用 \code{language=} 指定語言名稱。以 Python 程式的寫法為例：

\begin{codeblock}[language=latex]
\begin{lstlisting}[language=python]
def solve(n):
    return sum(range(n))
\end{lstlisting}
\end{codeblock}

如果這段程式需要標題與編號，在方括號裡加上 \code{caption=} 與 \code{label=}。需要從內文引用時，用 \code{\lstref} 引用，寫法和前一章的圖表相同：

\begin{codeblock}[language=latex]
\begin{lstlisting}[language=python,caption={遞迴版本},label=lst:fib]
def fib(n):
    return n if n < 2 else fib(n - 1) + fib(n - 2)
\end{lstlisting}
\end{codeblock}

注意到：程式碼預設不顯示標題與編號，下面也採用不加標題的寫法。需要標題時，可以依上面的範例選填 \code{caption}，並用 \code{\lstref{lst:fib}} 引用加上標籤的程式碼。標題預設放在下方；若要移到上方，請將 \code{layout.caption-position.code} 設為 \code{"top"}。

\begin{lstlisting}[language=python]
def fib(n):
    return n if n < 2 else fib(n - 1) + fib(n - 2)
\end{lstlisting}

當程式碼超過一頁時，會與表格一樣自動跨頁，不必手動拆成數個區塊——直接使用 \code{lstlisting} 的區塊，以及加上標題的區塊，都支援跨頁；有標題時，整段程式仍共用一個編號與標題。若要把外部檔案的內容直接放進論文，可以用 \code{\lstinputlisting[language=python]{path/to/file.py}}。

若只是在句子中提到指令或檔名，請用 \code{\code} 包住文字，例如

\begin{lstlisting}[language=latex]
\code{make release}
\end{lstlisting}

這種行內寫法會保留等寬字型，字級與基線則和本文一致，不會另外加上底色或框線。因為它是逐字排版，不能寫在 \code{\section} 或 \code{\caption} 這類的大括號引數裡；那些地方請改用 \code{\texttt}，需要顯示反斜線等特殊符號時，再加上 \code{\textbackslash}。

如果需要調整程式碼區塊外觀時，請打開 \code{tw-thesis.sty}，修改程式碼區塊那一節最上方的主題設定即可。此處的設定會套用到共用的程式碼樣式，不必逐一修改各章範例；其中 \code{\twcodesize} 控制區塊字級，行內程式碼仍會跟隨所在段落的字級。以下列出各項設定：

\begin{lstlisting}[language=latex]
\definecolor{twcodebg}{gray}{0.973}      % 底色
\definecolor{twcodeblue}{HTML}{4B69C6}   % 行內程式碼與關鍵字的顏色
\newcommand\twcodesize{\fontsize{11.4}{15}\selectfont}   % 區塊字級
\setmonofont{DejaVu Sans Mono}           % 等寬字體，可以放自己的字型
\end{lstlisting}

\begin{tlt}{程式碼樣式參數}{tbl:code-theme}{L{0.28\linewidth}L{0.62\linewidth}}
參數 & 用途與範例 \\
\midrule
\code{\setmonofont} & 等寬字型；例：\code{\setmonofont{DejaVu Sans Mono}}。 \\
\code{\twcodesize} & 區塊字級；例：本文的 0.95 倍，行內程式碼仍跟隨本文。 \\
\code{twcodebg} & 區塊底色；例：\code{\definecolor{twcodebg}{gray}{0.973}}。 \\
\code{twcodeblue} & 行內程式碼與關鍵字的顏色；例：\code{4B69C6}。 \\
\code{twcodegray} & 註解的顏色。 \\
\end{tlt}

\section{附錄}

需要收錄問卷、補充資料或較長的推導時，可以放到 \code{contents/back/} 的附錄檔案。請先在 \code{main.tex} 寫一次 \code{\appendix}，再填入標題與內文。以下示範第一份附錄的寫法，編譯後會使用 A、B 等字母編號：

\begin{lstlisting}[language=latex]
\chapter{問卷內容}\label{appendix-a}

\section{第一部分}
\end{lstlisting}

\code{main.tex} 裡的 \code{\appendix} 會切換編號方式，並把章的計數歸零。因此接下來的章標題會顯示為「附錄 A」，節標題則使用 \code{A.1}。替換內容時，只需修改「問卷內容」與各節標題，不必在標題中手動輸入字母或編號。

若要新增第二份附錄，請建立另一個 \code{.tex} 檔案，同樣用 \code{\chapter} 寫標題，但不要再寫一次 \code{\appendix}。接著到 \code{main.tex}，在第一份附錄後面加入新檔案的 \code{\input}，編號就會接續為 B。

\code{\appendix} 之後編號的規則是固定的，第二份、第三份附錄仍使用相同設定，不要自己改成 \code{B.1} 或 \code{C.1}。

\section{學校規範更新}

各校的論文格式規範偶爾會修訂，而模板裡的 \code{schools/*.yml} 和 \code{docs/} 收錄的規範 PDF，都只是某個時間點的版本。為了避免在不知情的情況下照著舊規範排版，搭配模板使用的管理工具 \code{tw-thesis} 會定時抓取是否有新的論文規範。

\subsection{提醒運作方式}

另一個專案 \href{https://github.com/TW-Thesis/tw-thesis-updater}{\texttt{tw-thesis-updater}} 每週自動抓取各校的官方規範文件，並公開一份清單，記錄每間學校目前最新版本的日期與雜湊值。在專案資料夾內執行 \code{tw-thesis} 開啟管理介面後，首頁會拿 \code{config.yml} 選定的學校，比對專案 \code{docs/sources.lock.json} 裡記錄的版本；兩者不同便會出現黃色提醒框，說明該校規範已修訂、附上新規範的下載連結。

\begin{lstlisting}[language=bash]
tw-thesis    # 在專案資料夾內執行；首頁會顯示提醒
\end{lstlisting}

該檢查只針對選擇的那所學校，一天最多連線一次；離線、或清單尚未發布時不會顯示任何錯誤，也不會影響編譯。

\subsection{提醒後續事宜}

程式僅負責提醒，不會自動修改任何排版設定，也不會判斷新規範和舊規範差在哪裡：各校、各系所對規範的解讀不盡相同，有些條文（例如封面格式、裝訂順序）尚要符合系所額外的規定，所以是否要修改、修改幅度仍由使用者決定。看到提醒後，建議依序處理：

\begin{enumerate}
\item 從提醒框的連結下載新規範，並實際從頭讀過一遍，切勿只看公告的標題。如果學校網站同時提供「修正對照表」或「修訂說明」，先看那份就能最快能找出有哪幾條規則修改了。
\item 對照 \code{schools/<學校代碼>.yml}（例如政大是 \code{schools/nccu.yml}），逐項確認邊界、字級、行距、標題格式、封面與審定書、頁碼、參考文獻格式等設定是否需要調整。各項目的設定方式見\ref{cha:layout}；若規範要求的項目在學校設定檔裡找不到對應欄位，先到 \code{schools/default.yml} 查有沒有同名的設定。
\item 想先試試看、不想直接動學校設定檔時，可以把要調整的項目寫進自己的 \code{config.yml}。它的優先順序高於學校設定，確認結果沒問題後，再決定要不要正式改 \code{schools/}。
\item 需要的話，把新的 PDF 放進 \code{docs/} 取代舊檔（檔名內含版本日期），以後想對照時才不會拿到舊版本。
\item 調整完成後重新編譯，逐頁核對封面、審定書、目次、正文第一頁與參考文獻頁，最好再請指導教授或系辦幫忙確認一次。
\end{enumerate}

\subsection{免責聲明}

本模板的格式設定僅供參考，實際仍請以學校公布的論文規範及系所要求為準。規定往往不只一份：學校公布的是通則，系所、學院或指導教授另有補充要求（例如封面、裝訂、字體或圖表格式）的情況並不少見，這些都不會出現在模板的設定裡。

此外，規範更新提醒是靠定期抓取學校網站，可能晚於學校實際公告才出現，也可能因為學校改版網站而暫時抓不到，所以沒有收到提醒，並不代表規範沒有修訂。提交論文前，請務必再向系所或圖書館確認最新要求。


\references                          % references

\appendix
\chapter{編譯問題排查}\label{appendix-a}

使用模板時，若遇到編譯失敗、圖片找不到或 PDF 沒有更新，可以先依本附錄的順序檢查。建議保留終端機中的錯誤訊息，從第一個錯誤指出的檔案與位置開始處理，修正後再重新編譯。

這份檔案也示範附錄的排版方式。換成自己的論文時，可以保留章標題的寫法，將標題與內文替換成問卷、補充數據或其他資料，再到 \code{main.tex} 確認附錄的裝訂順序。

\section{指令執行失敗}

若終端機顯示找不到 \code{xelatex} 或 \code{latexmk}，請先確認已完成安裝，再重新開啟終端機，執行

\begin{lstlisting}[language=bash]
xelatex --version
\end{lstlisting}

確認能夠顯示版本資訊後，再嘗試編譯。若找不到的是 \code{make}，則可以先使用以下指令，直接產生 PDF：

\begin{lstlisting}[language=bash]
texlua scripts/gen.lua
latexmk -xelatex main.tex
\end{lstlisting}

若錯誤訊息指出找不到 \code{main.tex} 或 \code{.tw-thesis-config.tex}，請檢查終端機目前所在的資料夾。編譯指令應在專案根目錄執行，也就是能看到 \code{main.tex}、\code{config.yml} 與 \code{contents/} 的位置；\code{.tw-thesis-config.tex} 則是 \code{scripts/gen.lua} 產生的，先執行過一次就會出現。若目前位於章節資料夾，請先切換回根目錄再試一次。

\section{修改設定後無法編譯}

修改設定後出現錯誤時，請先看訊息指出哪個檔案、哪一個欄位。若包含 \code{unknown key}，可以對照 \code{schools/default.yml} 檢查欄位拼字與所在層級；若是 \code{YAML} 解析錯誤，則先確認縮排使用空白，且冒號後保留空白，再檢查引號是否成對。

加入新設定時，請先找出檔案中是否已有相同區塊，避免在同一層重複建立 \code{layout:}。例如要同時修改字級與左側邊界，應把兩項設定放在同一個 \code{layout:} 下，並讓 \code{left} 再往內縮排一層：

\begin{lstlisting}[language=yaml]
layout:
  font-size: 12
  margin:
    left: 3.5
\end{lstlisting}

若一次修改了多個設定，暫時無法判斷錯誤來源，可以先將剛加入的內容複製到別處保存，再移除這些設定，確認原本的文件能正常編譯。之後逐項加回，每次存檔後檢查結果，較容易找到需要修正的位置。

\section{圖片或引用找不到}

圖片無法載入時，請先確認檔案確實放在 \code{figures/} 的指定位置，再比對路徑中的檔名、副檔名與大小寫。路徑是從專案根目錄算起的，不是從目前章節所在的資料夾開始。若曾經移動或重新命名圖片，也要一併修改章節中的路徑。

圖表引用顯示成 \code{??} 時，請檢查圖表是否已加上 \code{\label}，以及名稱是否與內文的 \code{\ref} 完全相同；也別忘了 \LaTeX{} 需要編譯兩次才會解出交叉引用，\code{latexmk} 會自動處理。若是文獻引用出錯，則到 \code{.bib} 檔案確認識別名稱，再檢查 \code{config.yml} 的 \code{bibliography.path} 是否包含該文獻檔。修正後重新編譯，確認引用已顯示為編號或文獻資料。

\section{PDF 沒有更新}

修改內容後，請先確認終端機是否已完成這次編譯，以及是否出現錯誤。編譯失敗時，資料夾中可能仍保留上一次的 \code{main.pdf}，因此能開啟 PDF 並不代表新內容已經寫入。請先修正錯誤，再重新執行編譯指令。

若使用 \code{make watch}，請確認監看程序仍在執行，並且修改的檔案已經存檔。編譯成功後畫面仍未更新時，可以關閉再重新開啟 PDF，同時確認檢視器開啟的是目前專案中的 \code{main.pdf}，避免看到先前複製到其他位置的版本。若編譯結果卡在舊的狀態，執行 \code{make clean} 清掉暫存檔後再編譯一次，通常就能解決。
     % appendix
\end{lstlisting}

\section{Code Blocks}

To list code in the thesis, create a code block right in the chapter file. The template applies a monospaced font and a light grey background and colours the syntax for the language you name. Choose the language name that fits, such as \code{python}, \code{bash} or \code{yaml}, and then enter the code.

A code block is made by wrapping the code to be shown in the \code{lstlisting} environment and naming the language with \code{language=} in square brackets. Here is how a Python program is written:

\begin{codeblock}[language=latex]
\begin{lstlisting}[language=python]
def solve(n):
    return sum(range(n))
\end{lstlisting}
\end{codeblock}

If the code needs a caption and a number, add \code{caption=} and \code{label=} in the square brackets. To refer to it from the text, use \code{\lstref}, written just like the figures and tables of the previous chapter:

\begin{codeblock}[language=latex]
\begin{lstlisting}[language=python,caption={Recursive version},label=lst:fib]
def fib(n):
    return n if n < 2 else fib(n - 1) + fib(n - 2)
\end{lstlisting}
\end{codeblock}

Note that code shows no caption or number by default, and the one below also uses the form without a caption. When you need a caption, add \code{caption} as in the example above, and refer to code that carries a label with \code{\lstref{lst:fib}}. The caption goes below by default; to move it above, set \code{layout.caption-position.code} to \code{"top"}.

\begin{lstlisting}[language=python]
def fib(n):
    return n if n < 2 else fib(n - 1) + fib(n - 2)
\end{lstlisting}

When code runs longer than a page, it breaks across pages like a table, with no need to split it into several blocks by hand: both a plain \code{lstlisting} block and one with a caption support page breaks, and with a caption the whole listing still shares one number and one caption. To put the contents of an external file into the thesis, use \code{\lstinputlisting[language=python]{path/to/file.py}}.

To mention a command or file name in a sentence, wrap the text in \code{\code}, for example

\begin{lstlisting}[language=latex]
\code{make release}
\end{lstlisting}

This inline form keeps the monospaced font, with the size and baseline matching the body, and adds no background or frame. Because it typesets its argument verbatim, it cannot be written inside the braces of commands such as \code{\section} or \code{\caption}; use \code{\texttt} there, and add \code{\textbackslash} when you need to show a backslash or other special symbols.

To change the look of code blocks, open \code{tw-thesis.sty} and edit the theme at the top of the code-block section. The settings apply to the shared code style, so you never change the examples in each chapter one by one; \code{\twcodesize} controls the block size, while inline code still follows the size of the paragraph it is in. The settings are listed below:

\begin{lstlisting}[language=latex]
\definecolor{twcodebg}{gray}{0.973}      % background
\definecolor{twcodeblue}{HTML}{4B69C6}   % inline code and keyword colour
\newcommand\twcodesize{\fontsize{11.4}{15}\selectfont}   % block size
\setmonofont{DejaVu Sans Mono}           % monospaced font, put your own here
\end{lstlisting}

\begin{tlt}{Code style parameters}{tbl:code-theme}{L{0.28\linewidth}L{0.62\linewidth}}
Parameter & Purpose and example \\
\midrule
\code{\setmonofont} & Monospaced font; e.g. \code{\setmonofont{DejaVu Sans Mono}}. \\
\code{\twcodesize} & Block size; e.g. 0.95 times the body, while inline code still follows the body. \\
\code{twcodebg} & Block background; e.g. \code{\definecolor{twcodebg}{gray}{0.973}}. \\
\code{twcodeblue} & Colour of inline code and keywords; e.g. \code{4B69C6}. \\
\code{twcodegray} & Colour of comments. \\
\end{tlt}

\section{Appendices}

To include a questionnaire, supplementary material or a long derivation, put it in an appendix file under \code{contents/back/}. Write \code{\appendix} once in \code{main.tex} and then enter the title and text. Below is how the first appendix is written; when compiled it is numbered with letters, A, B and so on:

\begin{lstlisting}[language=latex]
\chapter{Questionnaire}\label{appendix-a}

\section{Part One}
\end{lstlisting}

The \code{\appendix} in \code{main.tex} switches the numbering style and resets the chapter counter. The chapter headings that follow therefore show as ``Appendix A'' and section headings use \code{A.1}. When replacing the content, change only ``Questionnaire'' and the section headings; never type the letter or number into a heading by hand.

To add a second appendix, create another \code{.tex} file and write its title with \code{\chapter} as well, but do not write \code{\appendix} again. Then add the \code{\input} of the new file after the first appendix in \code{main.tex}, and the numbering continues as B.

The numbering rule after \code{\appendix} is fixed, so the second and third appendices use the same settings; do not change them to \code{B.1} or \code{C.1} yourself.

\section{Regulation Updates}

The universities' thesis format regulations are revised from time to time, and the \code{schools/*.yml} in the template and the regulation PDFs collected in \code{docs/} are only the version of one moment. To keep you from typesetting to an old regulation without knowing it, the management tool \code{tw-thesis} that goes with the template regularly checks whether there is a new thesis regulation.

\subsection{How the Reminder Works}

Another project, \href{https://github.com/TW-Thesis/tw-thesis-updater}{\texttt{tw-thesis-updater}}, fetches the official regulation documents of each school every week and publishes a list recording the date and hash of the latest version for every school. When you run \code{tw-thesis} inside the project folder to open the management interface, its home page compares the school chosen in \code{config.yml} with the version recorded in the project's \code{docs/sources.lock.json}; if they differ, a yellow notice appears saying that the school's regulation has been revised, with a download link to the new one.

\begin{lstlisting}[language=bash]
tw-thesis    # run inside the project folder; the home page shows the reminder
\end{lstlisting}

The check covers only the school you selected and connects at most once a day; when you are offline, or the list has not been published yet, no error is shown and compiling is not affected.

\subsection{What to Do After a Reminder}

The program only reminds you: it does not change any layout setting by itself, and it does not work out what differs between the new and old regulation. Schools and departments read the regulations differently, and some provisions (such as the cover format and binding order) must also meet the department's extra rules, so whether to change anything, and how much, is left to you. After seeing a reminder, work through the following in order:

\begin{enumerate}
\item Download the new regulation from the link in the notice and actually read it from the start; do not go by the headline of the announcement alone. If the school's site also offers a ``comparison of amendments'' or a ``revision note'', look at that first, since it is the fastest way to find which rules changed.
\item Compare it with \code{schools/<school code>.yml} (for National Chengchi University that is \code{schools/nccu.yml}) and check one by one whether the margins, font sizes, line spacing, heading format, cover and verification letter, page numbers and reference format need adjusting. See \ref{cha:layout} for how each item is set; if an item the regulation requires has no matching field in the school file, first look in \code{schools/default.yml} for a setting of the same name.
\item To try something out without touching the school file, write the items you want to change into your own \code{config.yml}. Its priority is higher than the school settings, and once the result looks right you can decide whether to change \code{schools/} for good.
\item If needed, put the new PDF in \code{docs/} to replace the old one (the file name carries the version date), so that you do not pick up the old version when you want to compare later.
\item After adjusting, compile again and check the cover, verification letter, table of contents, first page of the body and reference page one by one; it is best to have your advisor or the department office confirm once more.
\end{enumerate}

\subsection{Disclaimer}

The template's format settings are for reference only; in the end follow the thesis regulations published by your school and the requirements of your department. There is often more than one set of rules: what the school publishes is the general rule, and it is not unusual for departments, colleges or advisors to add requirements (for cover, binding, fonts or the format of figures and tables), none of which appear in the template's settings.

In addition, the regulation reminder works by fetching school websites on a schedule, so it may appear later than the school's actual announcement, and it may fail to fetch for a while when a school redesigns its site; not receiving a reminder therefore does not mean the regulation has not been revised. Before submitting your thesis, be sure to confirm the latest requirements once more with your department or the library.
